\documentclass[UTF8,a4paper,10pt]{ctexart}
\usepackage{../../../../../../common/ipara}
\usepackage{amsmath,amssymb,booktabs,tabularx,tikz}
\usetikzlibrary{arrows.meta,positioning}
\begin{document}
\section*{H-B02\quad 局部模型与区间定理：中值点、余项与误差控制}
\textbf{边界：}本桥只负责把局部近似升级为区间结论。Taylor 多项式与中值定理的完整机制仍由 Topic03、Topic04 持有。
\begin{tcolorbox}[title=母接口]
目标误差 $\longrightarrow$ 减去低阶模型 $\longrightarrow$ 构造零点函数
$\longrightarrow$ Rolle 链 $\longrightarrow$ 中值点余项/导数估计。
\end{tcolorbox}
\begin{center}
\begin{tikzpicture}[node distance=8mm, every node/.style={draw,rounded corners,align=center,minimum width=3.1cm,minimum height=8mm,font=\small}, >=Latex]
\node (a) {局部模型\\$P_n(x)$};
\node (b) [right=of a] {消去已知零点\\$\Phi=f-P_n-K\Pi$};
\node (c) [right=of b] {Rolle 链\\重复取导};
\node (d) [right=of c] {余项\\误差界};
\draw[->,thick] (a)--(b); \draw[->,thick] (b)--(c); \draw[->,thick] (c)--(d);
\end{tikzpicture}
\end{center}
\subsection*{1.\ Taylor 余项的接口}
若 $f\in C^{n+1}$，则在区间内存在 $\xi$ 使
\[
f(x)=\sum_{k=0}^{n}\frac{f^{(k)}(a)}{k!}(x-a)^k+
\frac{f^{(n+1)}(\xi)}{(n+1)!}(x-a)^{n+1}.
\]
因此若 $|f^{(n+1)}(t)|\le M$，立即得到
\[
|R_n(x)|\le \frac{M}{(n+1)!}|x-a|^{n+1}.
\]
这里的停止条件是：学生能说明“$\xi$ 只保证存在，$M$ 来自整个区间上界”，即可进入数值估计。
\subsection*{2.\ 插值余项的构造}
取节点 $x_0,\ldots,x_n$，令 $\Pi(x)=\prod_{i=0}^{n}(x-x_i)$。对
\[
\Phi(t)=f(t)-P(t)-K\Pi(t)
\]
选择 $K$ 使 $\Phi(x)=0$，并利用节点零点与 Rolle 定理逐层取导，得到
\[
f(x)-P(x)=\frac{f^{(n+1)}(\xi)}{(n+1)!}\Pi(x).
\]
误差控制仍先求 $\max|\Pi|$，再乘以导数上界。
\subsection*{3.\ 最小例子与调用}
对 $f(x)=e^x$ 在 $a=0$ 作一阶 Taylor，若 $|x|\le1$，存在 $\xi$ 介于 $0,x$ 之间使
\[
e^x=1+x+\frac{e^\xi}{2}x^2,\qquad
|e^x-1-x|\le\frac e2x^2.
\]
局部系数来自 Topic03，区间导数上界和中值点余项由本 Bridge 接通。题面要求误差界或“存在 $\xi$”时调用。

\subsection*{4.\ 合法性与反例边界}
\begin{itemize}
\item 中值点一般不能固定为端点、中心或某个“看起来合理”的点；
\item 多个中值点余项必须保证对应区间和正则性条件，不能把存在性写成同一点；
\item 只有局部导数信息时，不能推出整个区间的统一误差界。
\end{itemize}
\begin{tcolorbox}[title=检查句]
“我用的是哪条中值定理？零点在哪里？导数上界覆盖哪个区间？余项中的 $\xi$ 是否仍只是存在性变量？”
\end{tcolorbox}
\end{document}
